Building on the earlier Future Circular Collider (FCC) Conceptual Design Study conducted between 2014 and 2018, the FCC Feasibility Study (2021–2025) has been undertaken by a robust international collaboration, now comprising over 160 institutes worldwide.  The FCC “integrated programme”, developed in the framework of the Feasibility Study, consists of an initial electron-positron collider, the FCC-ee, which could be followed by a proton-proton collider, the FCC-hh. This staging takes into account the physics priorities as formulated in the updates of the European Strategy for Particle Physics of 2012 and 2020, as well as the relative technology readiness and costs of the FCC-ee and FCC-hh.

Over the years, I have closely followed the steady progress of the study, representing the FCC collaboration at the international steering committee and participating in annual FCC Week meetings, which include sessions of the International Collaboration Board. The commitment and enthusiasm of the members of the collaboration has always been impressive. The collective effort is clearly visible. Participation by students and early-career researchers is increasing. There is a shared determination and momentum to move forward.

The strong international collaboration around the FCC and its global network provide a solid foundation for the future of this project. The FCC community continues to grow, with increasing engagement from new institutes and partners worldwide. This broad support will be essential as the project enters its next phase.

The FCC Feasibility Study demonstrates not only the technical viability of the project, but also the strength of the international community that supports it. As we move towards the next step in the decision-making phase, this collective effort is key to showing a possible path forward. The FCC promises far-reaching scientific opportunities and long-term benefits for innovation, training, and global collaboration in science and technology.

\begin{flushright}
\textbf{Philippe Chomaz}\\
CEA, Chair of the FCC International Collaboration Board
\end{flushright}